% !TeX root = main.tex
\section{Datos}
\label{sec:datos}

\subsection{Procedencia y adquisición}

Se utilizan imágenes de resonancia magnética de la colección
\textit{Vestibular-Schwannoma-SEG}, publicada por Shapey y colaboradores en
The Cancer Imaging Archive~\cite{clark2013}. La colección comprende 242 pacientes adultos con
schwannoma vestibular unilateral esporádico tratados mediante radiocirugía
estereotáctica con Gamma Knife. Cada caso incluye imágenes ponderadas en T1
con contraste y en T2 de alta resolución, junto con estructuras segmentadas
y datos de planificación de radioterapia. Los datos se distribuyen bajo la
licencia de atribución CC BY 4.0~\cite{shapeydatos}.

Las imágenes originales se adquirieron con un resonador Siemens Avanto de
1,5~T. La adquisición T1 con contraste empleó una secuencia MPRAGE, con una
matriz de $512\times512$ píxeles y una resolución en el plano de
$0{,}4\times0{,}4$~mm. La adquisición T2 de alta resolución empleó una
secuencia 3D CISS, con matrices de $384\times384$ o $448\times448$ píxeles y
una resolución en el plano de aproximadamente $0{,}5\times0{,}5$~mm. El
espesor de corte fue de 1,0 a 1,5~mm. Las estructuras de referencia fueron
delimitadas por el equipo clínico mediante Leksell GammaPlan, empleando
las imágenes T1 con contraste y T2 de alta resolución. Estos parámetros describen
la adquisición clínica de la colección~\cite{shapey2021}.

\subsection{Selección y preparación de las imágenes}

El conjunto de trabajo comprende 30 imágenes bidimensionales, una por cada
caso identificado desde \texttt{VS-SEG-001} hasta \texttt{VS-SEG-030}. Según
la selección documentada en la presentación del proyecto~\cite{presentacion}, 13 cortes
corresponden a T1 con contraste y 17 a T2 de alta resolución. La
tabla~\ref{tab:composicion-datos} resume esta composición y las dimensiones
verificadas en los archivos. La distribución de tamaños concuerda con las
matrices de adquisición publicadas~\cite{shapey2021}.

\begin{table}[htbp]
\centering
\caption{Composición del conjunto de trabajo. La secuencia corresponde a la
clasificación documentada en la presentación; las dimensiones y cantidades
se verificaron en los archivos del repositorio~\cite{cuaderno}.}
\label{tab:composicion-datos}
\begin{tabular}{lrr}
\toprule
Secuencia & Dimensiones (píxeles) & Cantidad \\
\midrule
T1 con contraste & $512\times512$ & 13 \\
T2 de alta resolución & $448\times448$ & 16 \\
T2 de alta resolución & $384\times384$ & 1 \\
\midrule
Total & & 30 \\
\bottomrule
\end{tabular}
\end{table}

Las imágenes se extrajeron y convirtieron a PNG, tanto en su versión sin
superposición como en su versión con la anotación roja. El repositorio conserva las primeras en
\texttt{data/images/}, en escala de grises de 8 bits, y las segundas en
\texttt{data/contours/}, con canales rojo, verde y azul. Ambos archivos
comparten el identificador del caso. El análisis utiliza estos cortes
bidimensionales como unidades de procesamiento.

\subsection{Referencia de segmentación}

La implementación extrae una máscara binaria de la anotación
roja~\cite{cuaderno}. Para el píxel
$i$, sean $r_i$, $v_i$ y $b_i$ las intensidades de los canales rojo, verde y
azul, respectivamente. La etiqueta de referencia se define como
\begin{equation}
 g_i =
 \begin{cases}
  1, & r_i > v_i \text{ y } r_i > b_i,\\
  0, & \text{en otro caso}.
 \end{cases}
 \label{eq:referencia-roja}
\end{equation}
La referencia contiene entre 70 y 2748 píxeles positivos por imagen. En los
30 casos se verificó que la imagen y su anotación tienen las mismas
dimensiones y que todos los píxeles exteriores a la región roja coinciden
exactamente. Esta comprobación respalda la correspondencia espacial entre
los archivos locales. La máscara se utiliza para evaluar la segmentación;
no se suministra como etiqueta de entrada a GrabCut.

La máscara obtenida mediante la ecuación~\eqref{eq:referencia-roja}
permite comparar la segmentación calculada con la región marcada en la
anotación de cada imagen.
